\documentclass[11pt,a4paper]{article}
\usepackage[T1]{fontenc}
\usepackage[utf8]{inputenc}
\usepackage[french]{babel}
\usepackage{lmodern}
\usepackage{geometry}
\usepackage{amsmath,amssymb}
\usepackage{booktabs}
\usepackage{graphicx}
\usepackage{tabularx}
\graphicspath{{./}{../figures/}}
\usepackage{float}
\usepackage{enumitem}
\usepackage{hyperref}
% typo_francaise.tex
% Règles typographiques françaises (guide du dactylographe).
% Inclusion : % typo_francaise.tex
% Règles typographiques françaises (guide du dactylographe).
% Inclusion : % typo_francaise.tex
% Règles typographiques françaises (guide du dactylographe).
% Inclusion : \input{typo_francaise} dans le préambule.
% Pré-requis : babel chargé avec l'option [french] (gère déjà les espaces
% fines insécables avant : ; ! ? et autour de « »).

% --- Espace fine insécable automatique avant : ; ! ? -------------------
% Force babel-french à insérer une fine insécable avant la ponctuation double
% quelle que soit la saisie (mot:, mot :, mot   : donnent le même rendu).
\frenchsetup{AutoSpacePunctuation=true}

% --- Guillemets français : taper "texte" produit « texte » ---------------
\usepackage{csquotes}
\MakeOuterQuote{"}

% --- Nombres : 1234567 -> 1 234 567 ; décimales avec virgule -------------
% Usage : \num{1234567}, \num{3,14}, \SI{12}{\kilo\meter}, \SI{50}{\percent}
\usepackage{siunitx}
\sisetup{
  locale                  = FR,
  group-separator         = {\,},   % espace fine entre milliers
  group-minimum-digits    = 4,      % regrouper à partir de 4 chiffres
  output-decimal-marker   = {,},
  per-mode                = symbol,
  detect-all                        % suit la fonte du texte courant
}

% --- Micro-typographie : kerning, justification, débord en marge ---------
\usepackage{microtype}

% --- Tirets ---------------------------------------------------------------
% Sous LaTeX : --- = tiret cadratin (—), -- = demi-cadratin (–), - = trait d'union.
% Aucune configuration nécessaire, juste un rappel d'usage.

% --- Note d'usage ---------------------------------------------------------
% Espaces insécables manuelles (toujours nécessaires hors ponctuation forte) :
%   M.~Dupont, 12~km, n\textsuperscript{o}~3, p.~42, etc.
% Les espaces fines insécables avant : ; ! ? sont automatiques (babel french).
 dans le préambule.
% Pré-requis : babel chargé avec l'option [french] (gère déjà les espaces
% fines insécables avant : ; ! ? et autour de « »).

% --- Espace fine insécable automatique avant : ; ! ? -------------------
% Force babel-french à insérer une fine insécable avant la ponctuation double
% quelle que soit la saisie (mot:, mot :, mot   : donnent le même rendu).
\frenchsetup{AutoSpacePunctuation=true}

% --- Guillemets français : taper "texte" produit « texte » ---------------
\usepackage{csquotes}
\MakeOuterQuote{"}

% --- Nombres : 1234567 -> 1 234 567 ; décimales avec virgule -------------
% Usage : \num{1234567}, \num{3,14}, \SI{12}{\kilo\meter}, \SI{50}{\percent}
\usepackage{siunitx}
\sisetup{
  locale                  = FR,
  group-separator         = {\,},   % espace fine entre milliers
  group-minimum-digits    = 4,      % regrouper à partir de 4 chiffres
  output-decimal-marker   = {,},
  per-mode                = symbol,
  detect-all                        % suit la fonte du texte courant
}

% --- Micro-typographie : kerning, justification, débord en marge ---------
\usepackage{microtype}

% --- Tirets ---------------------------------------------------------------
% Sous LaTeX : --- = tiret cadratin (—), -- = demi-cadratin (–), - = trait d'union.
% Aucune configuration nécessaire, juste un rappel d'usage.

% --- Note d'usage ---------------------------------------------------------
% Espaces insécables manuelles (toujours nécessaires hors ponctuation forte) :
%   M.~Dupont, 12~km, n\textsuperscript{o}~3, p.~42, etc.
% Les espaces fines insécables avant : ; ! ? sont automatiques (babel french).
 dans le préambule.
% Pré-requis : babel chargé avec l'option [french] (gère déjà les espaces
% fines insécables avant : ; ! ? et autour de « »).

% --- Espace fine insécable automatique avant : ; ! ? -------------------
% Force babel-french à insérer une fine insécable avant la ponctuation double
% quelle que soit la saisie (mot:, mot :, mot   : donnent le même rendu).
\frenchsetup{AutoSpacePunctuation=true}

% --- Guillemets français : taper "texte" produit « texte » ---------------
\usepackage{csquotes}
\MakeOuterQuote{"}

% --- Nombres : 1234567 -> 1 234 567 ; décimales avec virgule -------------
% Usage : \num{1234567}, \num{3,14}, \SI{12}{\kilo\meter}, \SI{50}{\percent}
\usepackage{siunitx}
\sisetup{
  locale                  = FR,
  group-separator         = {\,},   % espace fine entre milliers
  group-minimum-digits    = 4,      % regrouper à partir de 4 chiffres
  output-decimal-marker   = {,},
  per-mode                = symbol,
  detect-all                        % suit la fonte du texte courant
}

% --- Micro-typographie : kerning, justification, débord en marge ---------
\usepackage{microtype}

% --- Tirets ---------------------------------------------------------------
% Sous LaTeX : --- = tiret cadratin (—), -- = demi-cadratin (–), - = trait d'union.
% Aucune configuration nécessaire, juste un rappel d'usage.

% --- Note d'usage ---------------------------------------------------------
% Espaces insécables manuelles (toujours nécessaires hors ponctuation forte) :
%   M.~Dupont, 12~km, n\textsuperscript{o}~3, p.~42, etc.
% Les espaces fines insécables avant : ; ! ? sont automatiques (babel french).

\geometry{margin=2.6cm}
\setlist[itemize]{leftmargin=*,nosep}
\setlist[enumerate]{leftmargin=*,nosep}

\newcommand{\modelrepo}{https://github.com/ajossto/LIED-StageAJS}

\newcommand{\monimage}[4][0.95]{%
  \begin{figure}[H]
    \centering
    \includegraphics[width=#1\linewidth]{#2}
    \caption{#3}
    \label{#4}
  \end{figure}%
}

\newcommand{\monimagelarge}[4][0.95]{%
  \begin{figure}[H]
    \centering
    \includegraphics[width=#1\textwidth]{#2}
    \caption{#3}
    \label{#4}
  \end{figure}%
}

\title{Familles de lois pour la distribution en taille et en revenu\\
de la population d'entités --- \texttt{modele-27-04-WIP}}
\author{Note technique --- Claude Code, à la demande d'Anatole Joseph-Stouls}
\date{2 juillet 2026}

\begin{document}
\maketitle

\begin{abstract}
On cherche la famille de lois (ou combinaison de lois) qui décrit le mieux, à
un pas de temps stationnaire, la distribution en taille des entités
(\texttt{actif\_total} / \texttt{passif\_total}) et la distribution de leurs
revenus (\texttt{revenu\_total} = extraction $\alpha\sqrt{P}$ + intérêts
reçus) dans les simulations du modèle \texttt{modele-27-04-WIP}. L'hypothèse
de départ --- un mélange de deux lois log-normales, jugé insatisfaisant ---
est testée contre une échelle de familles à 1 à 5 paramètres (exponentielle,
Fisk, Singh-Maddala, GB2, mélanges). Le mélange à 2 log-normales gagne
systématiquement au sens de l'AIC/BIC, sur 163 simulations stationnaires
regroupées par régime physique. Mais l'origine du mélange n'est pas ce que
l'hypothèse de départ supposait~: ses deux composantes correspondent à deux
\emph{cohortes d'âge} (entités fondatrices créées à $t=0$ contre entités
entrantes créées en continu), pas à deux types économiques distincts. Un
test de dérive avec la durée simulée montre en outre que la queue de la
distribution s'alourdit avec le temps~: ce n'est pas une loi d'échelle
stable, mais un artefact transitoire lié au protocole d'amorçage.
\end{abstract}

\tableofcontents

\section{Contexte et question posée}

Le modèle \texttt{modele-27-04-WIP} (source dans
\texttt{modele-27-04-WIP/src}) simule une population d'entités économiques
créées en continu, empruntant et prêtant un actif unique (le joule), avec
faillites possibles. La question posée~: quelle famille de lois décrit le
mieux, pour un nombre de paramètres donné (1 à 5), la distribution
transversale (à un instant donné, sur toutes les entités vivantes)
\begin{itemize}
  \item de la \textbf{taille} des entités (\texttt{actif\_total} ou
  \texttt{passif\_total})~;
  \item de leur \textbf{revenu} (\texttt{revenu\_total}, extraction
  $\Pi = \alpha\sqrt{P}$ plus intérêts reçus)~?
\end{itemize}

L'ajustement actuellement utilisé dans \texttt{analysis.py}
(\texttt{\_fit\_lognormal\_mixture}) est un mélange de deux lois
log-normales, jugé peu satisfaisant par l'auteur, qui soupçonnait deux
populations qui se chevauchent sans être certain de la nature de cette
hétérogénéité.

\section{Méthode}

\subsection{Sélection des simulations}

Sur les 805 simulations disponibles issues de ce modèle (13 lancées
directement, 792 issues de l'étude de sensibilité
\texttt{etude\_sensibilite\_27\_04\_wip}, même code source), \textbf{163
seulement sont stationnaires} au sens déjà calculé par le harnais de
sensibilité (\texttt{run\_simulation.py:detect\_regime\_diagnostics}) ---
queue bornée (\texttt{bounded\_tail}) et flux de faillites en équilibre avec
le flux de création (\texttt{flow\_balanced}). Ce filtrage est nécessaire~:
le processus brownien géométrique sur $\alpha$ documenté dans
\texttt{simulation.py} porte un \emph{drift positif} explicitement signalé
comme problème dans le code source, et la majorité des simulations ne
convergent jamais vers un régime stable.

\subsection{Principe~: pas de pooling}

Les valeurs ne sont \textbf{jamais regroupées} entre simulations ni entre
pas de temps avant ajustement --- regrouper des runs à paramètres différents,
ou des pas de temps différents d'une même trajectoire encore transitoire,
fabrique artificiellement de la multimodalité qui n'existe dans aucune
simulation individuelle. La méthode retenue est donc~: un ajustement par
paramètres maximum de vraisemblance (MLE) sur \textbf{un seul instantané
stationnaire par simulation}, puis agrégation des verdicts (famille
gagnante, valeurs de paramètres) sur l'ensemble des simulations d'un même
groupe.

Les simulations sont regroupées par les leviers physiques du modèle~:
\texttt{n\_candidats\_pool} ($k$, taille du pool de négociation crédit ---
$k\le 2$ : régime stable, $k\ge 3$ : régime critique auto-organisé selon
\texttt{config.py}), homogénéité de $\alpha$
($\alpha_{\min}=\alpha_{\max}$~?) et $\sigma_\alpha$ (volatilité du
brownien sur $\alpha$). Un regroupement par apprentissage non supervisé sur
la forme des distributions aurait été circulaire (on aurait regroupé les
runs sur l'objet même qu'on cherche à caractériser).

L'instant de mesure retenu par simulation est le pas recommandé par le
harnais (\texttt{t\_measure} / \texttt{tail\_start} / \texttt{t\_regime}
selon le schéma disponible), ou à défaut le dernier instantané disponible.

\subsection{Échelle de familles testées}

Treize familles ont été ajustées par MLE sur les valeurs strictement
positives, comparées par AIC et BIC~:

\begin{center}
\begin{tabularx}{\linewidth}{cX}
\toprule
Paramètres & Familles \\
\midrule
1 & Exponentielle, Pareto ($x_{\min}$ fixé au minimum observé) \\
2 & Log-normale, Gamma, Weibull, Fisk (log-logistique) \\
3 & Log-normale à 3 paramètres, Gamma généralisée, Singh-Maddala
    (Burr~XII), Dagum (Burr~III) \\
4 & GB2 (Generalized Beta of the Second Kind, McDonald 1984) --- ajustée
    par MLE numérique (implémentation propre, pas disponible dans
    \texttt{scipy}) \\
5 & Mélange de 2 log-normales (méthode existante du dépôt), mélange
    log-normale + queue de Pareto (méthode existante du dépôt) \\
\bottomrule
\end{tabularx}
\end{center}

Les familles à 2--4 paramètres forment une hiérarchie \emph{emboîtée}~: Fisk
$\subset$ Singh-Maddala $\subset$ GB2 (Fisk = GB2 avec $p=q=1$,
Singh-Maddala = GB2 avec $q=1$). C'est volontaire~: cela permet de savoir
si ajouter un paramètre à une même famille structurelle suffit, ou s'il
faut changer de structure (mélange).

Le code complet est dans \texttt{scripts/} (\texttt{families.py} pour
l'échelle de familles, \texttt{lib.py} pour le chargement des données,
\texttt{batch\_run.py} pour l'exécution en parallèle, \texttt{aggregate.py}
pour l'agrégation, \texttt{tail\_test.py} pour le test de queue). Les
résultats bruts par simulation sont dans \texttt{results/} (JSON).

\section{Résultat 1 --- l'origine du "mélange" est démographique, pas
économique}
\label{sec:cohortes}

L'hypothèse de départ était que le mélange reflète deux populations
d'entités économiquement distinctes. Les données ne le confirment pas~:
elles pointent vers une origine \textbf{démographique}.

Sur \texttt{entity\_final\_state.csv} (état final de toutes les entités,
avec \texttt{creation\_step}), on calcule l'âge de chaque entité vivante en
fin de simulation. La corrélation entre âge et $\log(\text{actif\_total})$
est de \textbf{0,88 à 0,95} selon les groupes --- y compris dans les
groupes à $\alpha$ \emph{hétérogène} ($\sigma_\alpha=0{,}02$), où l'âge
domine toujours le type de productivité comme facteur explicatif de la
taille.

Plus frappant~: la distribution des âges elle-même est nettement bimodale,
avec un \textbf{trou démographique} quasi total entre les deux modes
(figure~\ref{fig:cohorte})~: une entité créée après $t=0$ soit fait faillite
tôt (dans les premiers 10--15\,\% de la durée simulée), soit devient
quasiment immortelle. Il n'existe presque aucune entité d'âge intermédiaire.
Sur le groupe dominant ($k=4$, $\alpha$ homogène, $\sigma=0$, $n=54$
simulations), la fraction de "fondateurs" (âge~$>85\,\%$ de la durée
simulée) varie de 10 à 35\,\% selon la simulation, contre 0 à 11\,\%
d'entités d'âge intermédiaire.

\monimagelarge{age_cohorte.png}{Structure démographique de la population
vivante en fin de simulation (run représentatif, $k=4$, $\alpha$ homogène,
$\sigma=0$, $T=999$). À gauche~: taille en fonction de l'âge, échelle log
en ordonnée --- corrélation 0,88. À droite~: distribution des âges, avec le
trou démographique entre les entrants récents et les fondateurs
survivants.}{fig:cohorte}

Ce schéma s'explique par le protocole d'amorçage du modèle~: 100 entités
sont créées \emph{simultanément} à $t=0$
(\texttt{n\_entites\_initiales}), puis les entités suivantes arrivent en
continu (processus de type Poisson, taux \texttt{lambda\_creation}). Les
100 fondateurs, qui n'ont jamais été renouvelés, compoundent leur
croissance sur toute la durée de la simulation sans jamais mourir une fois
établis --- ils forment la queue lourde. Les entrants continus, eux,
meurent presque tous rapidement~; les rares survivants forment le corps
étroit de la distribution.

\textbf{Conclusion~:} l'intuition de deux populations qui se chevauchent
est correcte, mais ce sont deux \emph{cohortes d'âge} (fondateurs contre
entrants), pas deux classes économiques différentes. C'est un artefact du
protocole d'initialisation de la simulation, pas une loi structurelle de
l'économie modélisée.

\section{Résultat 2 --- hiérarchie des familles par nombre de paramètres}

Le tableau~\ref{tab:hierarchie} donne la famille gagnante (au sens de
l'AIC) pour chaque nombre de paramètres, sur les quatre groupes
comportant au moins 15 simulations stationnaires (tous à $\alpha$
homogène, $\sigma=0$)~; le résultat est identique pour
\texttt{actif\_total} et \texttt{revenu\_total}, et remarquablement stable
d'un groupe à l'autre.

\begin{table}[H]
\centering
\small
\begin{tabular}{ccllllc}
\toprule
$k$ & $n$ & 1p & 2p & 3p & 4p & 5p \\
\midrule
3 & 22 & Exponentielle & Fisk & Singh-Maddala & GB2 & Mélange log-log \\
4 & 54 & Exponentielle & Fisk & Singh-Maddala & GB2 & Mélange log-log \\
5 & 26 & Exponentielle & Fisk / log-nor. & Singh-Maddala & GB2 & Mélange log-log \\
6 & 15 & Exponentielle & Fisk / log-nor. & Singh-Maddala & GB2 & Mélange log-log \\
\bottomrule
\end{tabular}
\caption{Famille gagnante par nombre de paramètres, \texttt{actif\_total}
et \texttt{revenu\_total} (résultat identique), pour chaque valeur de $k=$
\texttt{n\_candidats\_pool} disposant d'au moins 15 simulations
stationnaires. "Mélange log-log" = mélange de 2 log-normales.}
\label{tab:hierarchie}
\end{table}

\textbf{Mise en garde sur le palier à 4 paramètres.} Le GB2 gagne l'AIC à
4 paramètres, mais son optimum se trouve en pratique sur une forme
\emph{dégénérée}~: le paramètre de forme $a$ vient buter sur la borne
numérique imposée ($a=50$), ce qui produit une fonction de survie qui
s'effondre presque verticalement (figure~\ref{fig:ccdf}, courbe verte
pointillée) --- un ajustement en quasi-fonction-marche, pas une forme
plausible. Autrement dit~: \emph{aucune famille lisse et unimodale ne
parvient à représenter un corps étroit accolé à une queue lourde}. C'est
précisément le signal qui justifie le saut vers un mélange explicite à 5
paramètres, cohérent avec la structure à deux cohortes du
résultat~\ref{sec:cohortes}. Le gain du mélange sur le GB2 est net
(médiane $\Delta\mathrm{BIC} \approx -100$ à $-120$ selon la variable, sur
le groupe $k=4$), donc ce n'est pas un artefact de sur-paramétrage~: c'est
une amélioration structurelle réelle.

\monimagelarge{ccdf_comparaison.png}{Fonction de survie empirique (points
noirs, échelle log-log) contre les quatre meilleures familles à 2, 3, 4 et
5 paramètres, pour \texttt{actif\_total} (gauche) et \texttt{revenu\_total}
(droite), sur une simulation représentative du groupe $k=4$, $\alpha$
homogène, $\sigma=0$.}{fig:ccdf}

\subsection{Paramètres typiques du mélange gagnant}

Sur le groupe dominant ($k=4$, $\alpha$ homogène, $\sigma=0$, $n=54$),
médiane des paramètres ajustés (convention \texttt{scipy.stats.lognorm}~:
$\mathrm{LN}(\mu,\sigma)$, échelle $=e^\mu$)~:

\begin{center}
\begin{tabular}{lcccc}
\toprule
Variable & $\lambda$ & Composante 1 (jeunes) & Composante 2 (fondateurs) & \\
\midrule
\texttt{actif\_total} & 0,51 & médiane 313 ($\sigma_1=0{,}060$) &
médiane 515 ($\sigma_2=0{,}440$) & \\
\texttt{revenu\_total} & 0,51 & médiane 17,6 ($\sigma_1=0{,}031$) &
médiane 23,7 ($\sigma_2=0{,}256$) & \\
\bottomrule
\end{tabular}
\end{center}

La composante~1 (entrants) est très resserrée ($\sigma_1 \approx
0{,}03$--$0{,}06$, coefficient de variation $\approx$ quelques \%)~: elle
correspond à une taille "typique" que les jeunes entités atteignent
rapidement puis autour de laquelle elles stagnent avant de faire faillite
ou d'être remplacées. La composante~2 (fondateurs) est beaucoup plus
dispersée ($\sigma_2 \approx 0{,}26$--$0{,}44$), cohérent avec une
croissance composée sur des durées très variables.

\section{Résultat 3 --- la queue n'est pas une loi de puissance stable}

La fonction de survie empirique décroche visiblement au-dessus de la
prédiction du mélange dans la queue extrême (figure~\ref{fig:ccdf})~: la
composante log-normale de queue décroît trop vite pour les plus grandes
valeurs observées. C'est probablement ce que l'auteur perçoit comme un
"mauvais ajustement" malgré un bon score AIC/BIC global (dominé par le
corps de la distribution, où l'accord est bon~: écart de Kolmogorov-Smirnov
médian de 0,04, sous le seuil critique à 5\,\% pour $n\approx 300$).

Un test de type Clauset--Shalizi--Newman (recherche du seuil $x_{\min}$ par
minimisation de la distance de Kolmogorov-Smirnov, puis test du rapport de
vraisemblance loi de puissance contre log-normale sur la queue au-delà de
$x_{\min}$) favorise la loi de puissance sur la log-normale dans la queue,
de façon quasi systématique (rapport de vraisemblance positif dans 95 à
100\,\% des simulations selon le groupe).

\textbf{Mais il ne s'agit probablement pas d'une vraie loi d'échelle.} Le
mécanisme identifié en section~\ref{sec:cohortes} --- une cohorte fixe de
fondateurs qui compound indéfiniment sans jamais être renouvelée ni
tuée --- est un mécanisme de \emph{diffusion log-normale}, pas le
mécanisme de renouvellement/rescaling (processus de Kesten) qui produit
usuellement les vraies lois de puissance en économie. Un test décisif~:
une loi d'échelle véritable a un exposant de queue stable quelle que soit
la durée simulée $T$~; un artefact transitoire dérive avec $T$.

\monimagelarge{derive_T.png}{Poids de la cohorte de fondateurs (gauche) et
exposant de queue estimé par la méthode Clauset--Shalizi--Newman (droite)
en fonction de la durée simulée $T$. Points gris~: simulations
individuelles~; losanges rouges~: médiane par palier de $T$.}{fig:derive}

Les deux quantités \textbf{dérivent nettement avec $T$}
(figure~\ref{fig:derive})~: la fraction de fondateurs passe de 34\,\% à
$T\approx 1000$ à 23\,\% à $T\approx 1500$ puis 13\,\% à $T\approx 3000$~;
l'exposant de queue passe de 5,4 à 5,0 puis 3,1 sur les mêmes paliers (plus
l'exposant est petit, plus la queue est lourde). Une vraie loi d'échelle
serait stable~; ici, la queue s'alourdit mécaniquement à mesure que les
fondateurs ont plus de temps pour composer leur croissance.

\textbf{Conclusion~:} le constat honnête n'est pas "la queue suit une loi
de puissance d'exposant $\alpha\approx5$", mais \emph{"la queue empirique
est systématiquement plus lourde que ne le prédit la composante
log-normale du mélange, et une queue de type Pareto/Singh-Maddala
l'approxime mieux localement --- sans que cela constitue une loi d'échelle
stationnaire du système"}.

\section{Recommandations pratiques}

\begin{itemize}
  \item Pour un usage descriptif courant (comparaison de scénarios,
  résumé d'une distribution), le \textbf{mélange à 2 log-normales} reste
  le meilleur choix "clé en main"~: il gagne l'AIC et le BIC dans la
  quasi-totalité des simulations stationnaires testées, pour les deux
  variables.
  \item \textbf{Ne pas interpréter ses deux composantes comme deux types
  économiques.} Elles correspondent à des cohortes d'âge, directement
  vérifiables via \texttt{creation\_step} dans
  \texttt{entity\_final\_state.csv}. Si une lecture par cohorte est
  utile, il vaut mieux conditionner explicitement sur l'âge plutôt que de
  laisser le mélange statistique jouer ce rôle par proxy.
  \item \textbf{Ne pas extrapoler la queue du mélange} pour des usages
  sensibles aux événements rares (risque systémique, taille de cascade
  maximale)~: la composante log-normale de queue sous-estime
  probablement les événements extrêmes. Un mélange log-normale + Pareto
  (déjà implémenté dans le dépôt sous
  \texttt{\_fit\_lognormal\_pareto}) est plus prudent pour ce type
  d'usage, même s'il perd légèrement au score AIC/BIC global.
  \item Le GB2 (4 paramètres) ne doit \textbf{pas} être utilisé tel quel
  malgré son bon AIC~: son optimum numérique est dégénéré sur ces
  données. Il a valeur de diagnostic (aucune famille lisse ne suffit),
  pas de recommandation.
\end{itemize}

\section{Limites et pistes}

\begin{itemize}
  \item L'essentiel des 163 simulations stationnaires est à $\alpha$
  \emph{homogène} ($\alpha_{\min}=\alpha_{\max}$)~; dans ce régime,
  \texttt{revenu\_total} est presque une fonction déterministe de
  \texttt{actif\_total} ($\Pi=\alpha\sqrt{P}$ à $\alpha$ constant), donc
  l'accord entre les résultats "taille" et "revenu" n'est pas une
  confirmation indépendante. Les groupes à $\sigma_\alpha>0$ seraient le
  vrai test d'indépendance, mais trop peu de simulations y sont
  stationnaires ($n\le 7$) pour trancher solidement.
  \item Le drift positif documenté du brownien sur $\alpha$
  (\texttt{simulation.py}, commentaire \textit{"GROS PROBLEME !!! DRIFT
  POSITIF"}) explique pourquoi 80\,\% des simulations disponibles ne sont
  pas stationnaires. Le corriger permettrait d'obtenir davantage de
  simulations stationnaires à $\sigma_\alpha>0$, et donc un vrai test
  indépendant taille/revenu.
  \item Piste directe pour vérifier le résultat~\ref{sec:cohortes}~: fitter
  séparément les deux sous-populations obtenues en conditionnant sur
  l'âge (via \texttt{creation\_step}) et vérifier que la bimodalité
  disparaît à l'intérieur de chaque cohorte.
  \item L'échelle de familles testées est positive uniquement.
  \texttt{revenu\_net} (revenu moins charges d'intérêts) peut être
  négatif et n'a pas été traité ici~; il demanderait une famille adaptée
  (normale asymétrique, ou décomposition explicite revenu brut moins
  charges).
\end{itemize}

\section*{Reproductibilité}

Le code d'ajustement (\texttt{scripts/families.py}), le chargement des
données (\texttt{scripts/lib.py}), l'exécution parallèle
(\texttt{scripts/batch\_run.py}), l'agrégation
(\texttt{scripts/aggregate.py}) et le test de queue
(\texttt{scripts/tail\_test.py}) sont fournis avec ce rapport, ainsi que
les résultats bruts par simulation (\texttt{results/*.json}, 163
fichiers). Dépôt~: \url{\modelrepo}.

\end{document}
